% ============================================================
%  Übungsaufgaben Template (my_dhbw Stil, uebung_*.tex)
%  Header "Übungsblatt", grün eingefärbte <details>-Lösungen.
%  Renderer: pdflatex (zweimal)
% ============================================================
\documentclass[11pt, a4paper]{article}

\usepackage[ngerman]{babel}
\usepackage{fontspec}
\setmainfont[
  Path=/usr/share/fonts/TTF/,
  Extension=.ttf,
  BoldFont=DejaVuSans-Bold,
  ItalicFont=DejaVuSans-Oblique,
  BoldItalicFont=DejaVuSans-BoldOblique
]{DejaVuSans}
\setsansfont[
  Path=/usr/share/fonts/TTF/,
  Extension=.ttf,
  BoldFont=DejaVuSans-Bold,
  ItalicFont=DejaVuSans-Oblique,
  BoldItalicFont=DejaVuSans-BoldOblique
]{DejaVuSans}
\setmonofont[Path=/usr/share/fonts/TTF/,Extension=.ttf]{DejaVuSansMono}
\usepackage[top=2cm, bottom=2cm, left=2.4cm, right=2.4cm, footskip=1cm]{geometry}
\usepackage{amsmath, amssymb}
\usepackage{xcolor}
\usepackage{enumitem}
\usepackage{fancyhdr}
\usepackage{lastpage}
\usepackage{tabularx}
\usepackage{booktabs}
\usepackage{microtype}
\usepackage{parskip}

\usepackage{array}
\usepackage{longtable}
\usepackage{ltablex}
\keepXColumns
\usepackage{colortbl}
\usepackage[most,breakable]{tcolorbox}
\usepackage{listings}
\usepackage[normalem]{ulem}
\usepackage{pdflscape}
\usepackage{titlesec}
\usepackage[colorlinks=true, linkcolor=accent, urlcolor=accent]{hyperref}

\renewcommand{\familydefault}{\sfdefault}

% ── Theme ─────────────────────────────────────────────────────
\definecolor{accent}{RGB}{0, 60, 113}
\definecolor{primaryblue}{RGB}{0, 60, 113}
\definecolor{loesung}{RGB}{0, 100, 0}
\definecolor{codebg}{RGB}{245, 245, 245}
\definecolor{figurebg}{RGB}{232, 241, 255}
\definecolor{tablehintbg}{RGB}{240, 255, 240}
\definecolor{solutionbg}{RGB}{240, 252, 240}
\definecolor{rulegray}{RGB}{180, 180, 180}
\definecolor{tableheadbg}{RGB}{0, 60, 113}

% ── Header/Footer ─────────────────────────────────────────────
\pagestyle{fancy}
\fancyhf{}
\renewcommand{\headrulewidth}{0.4pt}
\fancyhead[L]{\footnotesize\color{accent}\nichttitle}
\fancyhead[R]{\footnotesize\color{accent}\nichtsubtitle}
\fancyfoot[R]{\footnotesize\color{accent}Blatt \thepage\,/\,\pageref{LastPage}}

% ── Typografie ────────────────────────────────────────────────
\titleformat{\section}
    {\color{accent}\large\bfseries}{}{0pt}{}
    [\vspace{-4pt}\color{accent}\rule{\linewidth}{1.5pt}]
\titleformat{\subsection}
    {\normalsize\bfseries\color{accent!85!black}}{}{0pt}{}{}
\titleformat{\subsubsection}
    {\small\bfseries\color{accent!70!black}}{}{0pt}{}{}
\titlespacing*{\section}{0pt}{10pt}{3pt}
\titlespacing*{\subsection}{0pt}{6pt}{2pt}
\titlespacing*{\subsubsection}{0pt}{4pt}{1pt}

\setlength{\parindent}{0pt}
\setlength{\parskip}{6pt}
\renewcommand{\arraystretch}{1.2}
\setlist[itemize]{noitemsep, topsep=3pt, parsep=0pt, leftmargin=1.4em}
\setlist[enumerate]{noitemsep, topsep=3pt, parsep=0pt, leftmargin=1.6em}

% ── Boxen: solutionbox = Lösungsbox in grün ──────────────────
\newtcolorbox{figurebox}[1]{
  colback=figurebg, colframe=accent!50,
  title={Abbildung: #1}, fonttitle=\small\bfseries,
  breakable, before skip=6pt, after skip=6pt
}
\newtcolorbox{tablehintbox}[1]{
  colback=tablehintbg, colframe=green!50!black!50,
  title={Tabelle: #1}, fonttitle=\small\bfseries,
  breakable, before skip=6pt, after skip=6pt
}
\newtcolorbox{solutionbox}[1]{
  enhanced,
  colback=solutionbg, colframe=loesung,
  title={#1}, fonttitle=\small\bfseries,
  coltitle=white, colbacktitle=loesung,
  breakable, before skip=8pt, after skip=8pt
}

\lstset{
  basicstyle=\ttfamily\small,
  backgroundcolor=\color{codebg},
  breaklines=true,
  frame=single, framerule=0pt,
  xleftmargin=4pt, xrightmargin=4pt,
  aboveskip=6pt, belowskip=6pt,
  keepspaces=true
}

\providecommand{\nichttitle}{Übungsblatt}
\providecommand{\nichtsubtitle}{}

\renewcommand{\nichttitle}{Analysis 2}
\renewcommand{\nichtsubtitle}{Übungsblatt 4}


\begin{document}

\begin{center}
{\Large\bfseries\color{accent}\nichttitle}
\ifx\nichtsubtitle\empty\else
  \\[2pt]{\normalsize\color{accent!80!black}\nichtsubtitle}
\fi
\\[2pt]
\color{accent}\rule{0.4\linewidth}{0.8pt}\color{black}
\end{center}
\vspace{4pt}

% ── BODY ──────────────────────────────────────────────────────

\section{Übungsblatt 4 — Trigonometrie}


\noindent\textcolor{rulegray}{\rule{\linewidth}{0.4pt}}


\paragraph{Aufgabe 1 — Bogenmaß und Grundidentität \textit{(10 Punkte)}}\mbox{}\\[2pt]

a) Erkläre, warum bei der Definition trigonometrischer Funktionen üblicherweise das Bogenmaß (rad) und nicht das Gradmaß verwendet wird. Welche Konsequenz hat dies konkret für die Ableitungsregel von sin(t)? \textit{(4 P)}


b) Ein Winkel beträgt 75°. Rechne ihn ins Bogenmaß um (Bruch von π) und berechne damit cos²(75°) + sin²(75°). \textit{(3 P)}


c) Gegeben sei sin(t) = 3/5 mit t ∈ (0, π/2). Bestimme cos(t), tan(t) und sec(t) — ausschließlich mit Hilfe der Grundidentität, ohne arcsin zu verwenden. \textit{(3 P)}


\textbf{Lösung:}


a) Im Bogenmaß ist die Bogenlänge auf dem Einheitskreis gleich dem Winkel. Dadurch gilt der fundamentale Grenzwert lim\_\{t→0\} sin(t)/t = 1 nur im Bogenmaß. Konsequenz: Die elegante Ableitungsregel ∂/∂t sin(t) = cos(t) gilt nur im Bogenmaß. Bei Gradmaß würde ein Faktor π/180 als Kettenregel-Term auftreten, sodass ∂/∂t° sin(t°) = (π/180)·cos(t°).


b) 75° = 75·(π/180) = \textbf{5π/12 rad}.

Nach der Grundidentität gilt für jedes t: cos²(t) + sin²(t) = 1, also auch cos²(75°) + sin²(75°) = \textbf{1}.


c) Aus cos²(t) = 1 − sin²(t) = 1 − 9/25 = 16/25.

Da t ∈ (0, π/2): cos(t) > 0, also \textbf{cos(t) = 4/5}.

\textbf{tan(t) = sin(t)/cos(t) = (3/5)/(4/5) = 3/4}.

\textbf{sec(t) = 1/cos(t) = 5/4}.



\noindent\textcolor{rulegray}{\rule{\linewidth}{0.4pt}}


\paragraph{Aufgabe 2 — Wellenparameter und Ableitungs-Zyklus \textit{(15 Punkte)}}\mbox{}\\[2pt]

Eine harmonische Schwingung wird beschrieben durch

y(t) = 2.5 · sin(ωt) mit Frequenz f = 50 Hz.


a) Bestimme ω, λ und A. Schreibe y(t) explizit hin. \textit{(3 P)}


b) Berechne die ersten vier Ableitungen y′(t), y″(t), y‴(t), y⁽⁴⁾(t) und zeige damit, dass nach vier Ableitungen wieder die Ausgangsfunktion (bis auf einen Vorfaktor) entsteht. Welcher Vorfaktor ergibt sich? \textit{(6 P)}


c) Bei welchen Zeitpunkten t ∈ [0, 1/50 s] erreicht y(t) seine maximale Auslenkung, bei welchen ist die Geschwindigkeit y′(t) maximal? Begründe geometrisch über die Phasenverschiebung cos(t) = sin(t + π/2). \textit{(6 P)}


\textbf{Lösung:}


a) ω = 2πf = 2π·50 = \textbf{100π rad/s}.

λ = 2π/ω = 2π/(100π) = \textbf{1/50 = 0.02} (Einheit: s, da hier f in Hz, also λ als Periode T).

A = \textbf{2.5}.

y(t) = 2.5 · sin(100π·t).


b) Mit dem Zyklus sin → cos → −sin → −cos → sin und jedem Schritt zusätzlichem Faktor ω = 100π:

\begin{itemize}
  \item y′(t) = 2.5·100π·cos(100π·t) = 250π·cos(100π·t)
  \item y″(t) = −2.5·(100π)²·sin(100π·t) = −25000π²·sin(100π·t)
  \item y‴(t) = −2.5·(100π)³·cos(100π·t) = −2.5·10⁶π³·cos(100π·t)
  \item y⁽⁴⁾(t) = 2.5·(100π)⁴·sin(100π·t) = (100π)⁴ · y(t)
\end{itemize}

Nach vier Ableitungen: Vorfaktor \textbf{ω⁴ = (100π)⁴ ≈ 9.74·10⁸}, Funktionsform identisch.


c) Maximum von y(t): sin(100π·t) = 1 ⇔ 100π·t = π/2 + 2kπ ⇔ t = 1/200 + k/50.

In [0, 1/50] also \textbf{t = 1/200 s = 5 ms}.


Maximum von y′(t): cos(100π·t) = 1 ⇔ 100π·t = 0 + 2kπ ⇔ t = 0, 1/50.

In [0, 1/50] also \textbf{t = 0 und t = 1/50 s}.


Geometrisch: Da cos(x) = sin(x + π/2), läuft y′ der Funktion y um eine Viertelperiode (T/4 = 5 ms) \textbf{voraus}. Wenn die Auslenkung maximal ist, ist die Geschwindigkeit gerade Null — und umgekehrt; das ist genau die Phasenverschiebung um π/2.



\noindent\textcolor{rulegray}{\rule{\linewidth}{0.4pt}}


\paragraph{Aufgabe 3 — Partielle Ableitungen mit trig. Funktionen \textit{(20 Punkte)}}\mbox{}\\[2pt]

Gegeben sei f(x, y) = x · y · $e^{sin(x)}$ + sin²(x)·cos(y).


a) Berechne ∂f/∂x. Halte y konstant und führe die Produkt- und Kettenregel sauber aus. \textit{(8 P)}


b) Berechne ∂f/∂y. \textit{(4 P)}


c) Werte ∂f/∂x an der Stelle (x, y) = (0, 2) aus. \textit{(3 P)}


d) Sei zusätzlich g(x) = sin(arctan(x)). Zeige durch direkte Rechnung mit der Grundidentität (nicht mit der Tabelle), dass g′(x) = 1/√(x²+1) gilt. \textit{(5 P)}


\textbf{Lösung:}


a) Zwei Summanden, jeder mit Produkt-/Kettenregel:


Summand 1: ∂/∂x [x·y·$e^{sin(x)}$] = y·$e^{sin(x)}$ + x·y·cos(x)·$e^{sin(x)}$

Summand 2: ∂/∂x [sin²(x)·cos(y)] = 2·sin(x)·cos(x)·cos(y)


⇒ \textbf{∂f/∂x = y·$e^{sin(x)}$ + x·y·cos(x)·$e^{sin(x)}$ + 2·sin(x)·cos(x)·cos(y)}


b) Erster Summand: ∂/∂y [x·y·$e^{sin(x)}$] = x·$e^{sin(x)}$.

Zweiter Summand: ∂/∂y [sin²(x)·cos(y)] = −sin²(x)·sin(y).


⇒ \textbf{∂f/∂y = x·$e^{sin(x)}$ − sin²(x)·sin(y)}


c) Bei (0, 2): sin(0) = 0, cos(0) = 1, $e^{sin(0)}$ = e⁰ = 1.

∂f/∂x|\_\{(0,2)\} = 2·1 + 0·2·1·1 + 2·0·1·cos(2) = \textbf{2}.


d) Sei θ = arctan(x), also tan(θ) = x mit θ ∈ (−π/2, π/2). Dann ist cos(θ) > 0.

Aus 1 + tan²(θ) = sec²(θ) = 1/cos²(θ) folgt cos(θ) = 1/√(1+x²).

Damit sin(θ) = tan(θ)·cos(θ) = x/√(1+x²), also g(x) = x/√(1+x²).


Ableiten mit Quotientenregel/Produktregel:

g′(x) = [√(1+x²) − x · x/√(1+x²)] / (1+x²) = [(1+x²) − x²] / (1+x²)\textasciicircum{}\{3/2\} = 1/(1+x²)\textasciicircum{}\{3/2\}.


Hmm — Korrektur: das Ergebnis muss 1/√(x²+1) sein. Prüfung: Mit der Tabelle ist ∂/∂x[arctan(x)] = 1/(x²+1) und ∂/∂u[sin(u)] = cos(u). Kettenregel: g′(x) = cos(arctan(x)) · 1/(x²+1) = (1/√(1+x²)) · 1/(1+x²) = 1/(1+x²)\textasciicircum{}\{3/2\}.


Die Tabellenangabe „sin(arctan(x)) → 1/√(x²+1)" bezieht sich also auf den \textbf{Funktionswert} sin(arctan(x)) = x/√(x²+1) und dessen Vereinfachung — nicht auf eine Ableitung. Tatsächliche Ableitung:


⇒ \textbf{g′(x) = 1/(1+x²)\textasciicircum{}\{3/2\}}.


(Hinweis: Die Aufgabe stellt klar, dass die Tabelleneinträge sorgfältig auf Funktion vs. Ableitung zu lesen sind.)



\noindent\textcolor{rulegray}{\rule{\linewidth}{0.4pt}}


\paragraph{Aufgabe 4 — Transfer: Hyperbolische Funktionen in einem Neuron \textit{(15 Punkte)}}\mbox{}\\[2pt]

In einem neuronalen Netz wird tanh als Aktivierungsfunktion eines Neurons verwendet:

a(z) = tanh(z), z = w·x + b, mit w = 0.5, b = −1, Eingabe x = 4.


a) Begründe mit der Identität cosh²(t) − sinh²(t) = 1, dass tanh(t) ∈ (−1, 1) für alle t ∈ ℝ. \textit{(4 P)}


b) Zeige, dass tanh′(t) = 1 − tanh²(t). \textit{(5 P)}


c) Berechne die Ableitung der Aktivierung nach dem Gewicht w an der konkreten Stelle (x=4, w=0.5, b=−1). Nutze tanh(1) ≈ 0.7616. \textit{(6 P)}


\textbf{Lösung:}


a) Aus cosh²(t) − sinh²(t) = 1 folgt cosh²(t) = 1 + sinh²(t), also |cosh(t)| > |sinh(t)| für alle t (denn 1 + sinh² > sinh²). Da cosh(t) ≥ 1 > 0, folgt |sinh(t)/cosh(t)| < 1, also \textbf{tanh(t) ∈ (−1, 1)}.


b) tanh(t) = sinh(t)/cosh(t). Quotientenregel:

tanh′(t) = [cosh(t)·cosh(t) − sinh(t)·sinh(t)] / cosh²(t) = (cosh² − sinh²)/cosh² = 1/cosh²(t).


Nun 1/cosh²(t) = (cosh² − sinh²)/cosh² = 1 − sinh²/cosh² = \textbf{1 − tanh²(t)}. ✓


c) z = 0.5·4 + (−1) = 1.

a(z) = tanh(z), und ∂z/∂w = x = 4.

Kettenregel: ∂a/∂w = tanh′(z) · ∂z/∂w = (1 − tanh²(1)) · 4.


Mit tanh(1) ≈ 0.7616: tanh²(1) ≈ 0.5800.

1 − 0.5800 = 0.4200.

∂a/∂w ≈ 0.4200 · 4 = \textbf{1.680}.



\noindent\textcolor{rulegray}{\rule{\linewidth}{0.4pt}}


\paragraph{Aufgabe 5 — Fehler finden in einer Ableitungs-Tabelle \textit{(15 Punkte)}}\mbox{}\\[2pt]

Ein Kommilitone gibt dir folgende Tabelle (alle Ableitungen nach x):



{\small
\begin{tabularx}{\linewidth}{XXX}
\hline
\rowcolor{tableheadbg}
{\color{white}\textbf{Nr.}} & {\color{white}\textbf{f(x)}} & {\color{white}\textbf{behauptete Ableitung f′(x)}} \\[2pt]
\hline
\endhead
\hline
\endfoot
(i) & arcsin(x) & 1/√(1+x²) \\
(ii) & tan(x) & 1 + tan²(x) \\
(iii) & csc(x) & cot(x)·csc(x) \\
(iv) & √(1−cos²(x)), x∈(0,π) & cos(x) \\
(v) & sin(x)·cos(x) & cos²(x) − sin²(x) \\
\end{tabularx}
}


a) Identifiziere alle fehlerhaften Einträge und gib für jeden die korrekte Ableitung an. \textit{(8 P)}


b) Korrigiere den Fehler bei (i) und begründe, warum dieselbe Form (mit anderem Vorzeichen) für arccos(x) auftaucht. \textit{(4 P)}


c) Zeige für (iv), dass die behauptete Ableitung nicht nur falsch ist, sondern dass die Funktion sogar gleich sin(x) ist (im angegebenen Intervall). \textit{(3 P)}


\textbf{Lösung:}


a) Prüfung jeder Zeile:

\begin{itemize}
  \item (i) arcsin(x)′ = \textbf{1/√(1−x²)} (nicht 1+x²). \textbf{FEHLER}.
  \item (ii) tan(x)′ = sec²(x) = 1 + tan²(x). \textbf{KORREKT}.
  \item (iii) csc(x)′ = \textbf{−cot(x)·csc(x)} (Vorzeichen fehlt!). \textbf{FEHLER}.
  \item (iv) Im Intervall (0,π) gilt √(1−cos²(x)) = sin(x), also Ableitung = cos(x). \textbf{KORREKT} (Endergebnis), auch wenn der Weg über die Wurzel-Kettenregel umständlicher wäre.
  \item (v) (sin·cos)′ = cos·cos + sin·(−sin) = cos²(x) − sin²(x). \textbf{KORREKT}.
\end{itemize}

Fehlerhafte Einträge: \textbf{(i) und (iii)}.


b) Korrekt: arcsin′(x) = 1/√(1−x²). Da arcsin(x) + arccos(x) = π/2 (konstant), folgt durch Ableiten arccos′(x) = −arcsin′(x) = −1/√(1−x²). Daher dieselbe Form mit umgekehrtem Vorzeichen.


c) Nach der bewiesenen Identität gilt √(1−cos²(x)) = sin(x) genau für x ∈ (0, π) (denn dort ist sin(x) ≥ 0, sodass die positive Wurzel sin(x) selbst ergibt). Also ist die Funktion in diesem Intervall identisch zu sin(x), und ihre Ableitung trivialerweise cos(x). Die behauptete Ableitung ist also korrekt — der „Fehler", den die Tabelle suggerieren könnte (etwa eine komplizierte Wurzel), wäre, die Identität nicht zu erkennen.



\noindent\textcolor{rulegray}{\rule{\linewidth}{0.4pt}}


\paragraph{Aufgabe 6 — Vergleich: Trigonometrische vs. Hyperbolische Identität \textit{(15 Punkte)}}\mbox{}\\[2pt]

a) Stelle die beiden Pythagoras-artigen Identitäten gegenüber:

cos²(t) + sin²(t) = 1 (Kreis) gegen cosh²(t) − sinh²(t) = 1 (Hyperbel).

Beschreibe geometrisch, welche Kurve jeweils parametrisiert wird, und welche grundsätzliche Konsequenz das Vorzeichen für die Wertebereiche von sin/cos vs. sinh/cosh hat. \textit{(6 P)}


b) In der Fourier-Analyse tritt $e^{iωt}$ = cos(ωt) + i·sin(ωt) auf. Begründe damit, warum |$e^{iωt}$| = 1 für alle t — und warum die analoge Konstruktion mit sinh, cosh \textbf{nicht} auf den Einheitskreis führt. \textit{(5 P)}


c) Treffe eine begründete Entscheidung: Welche der beiden Funktionen — tanh oder eine reine sin-basierte Aktivierung — ist als neuronale Aktivierung sinnvoller? Nenne mindestens zwei Trade-offs, die die Wahl beeinflussen. \textit{(4 P)}


\textbf{Lösung:}


a) \textbf{Kreis}: (cos t, sin t) parametrisiert den Einheitskreis x²+y²=1, beschränkt. Folge: cos, sin sind beide auf [−1, 1] beschränkt, periodisch mit Periode 2π.

\textbf{Hyperbel}: (cosh t, sinh t) parametrisiert den \textbf{rechten Ast} der Hyperbel x²−y²=1. Folge: cosh(t) ≥ 1 (unbeschränkt nach oben), sinh(t) ∈ ℝ (unbeschränkt in beide Richtungen), nicht periodisch. Das Minuszeichen in cosh²−sinh²=1 spiegelt wider, dass die Hyperbel offen ist statt geschlossen.


b) |$e^{iωt}$|² = cos²(ωt) + sin²(ωt) = 1, also |$e^{iωt}$| = 1 für alle t (Punkt auf dem Einheitskreis in der komplexen Ebene). Eine analoge Konstruktion cosh(ωt) + i·sinh(ωt) hätte Betragsquadrat cosh²(ωt) + sinh²(ωt) = 1 + 2·sinh²(ωt) ≠ 1 (wächst exponentiell). Es ist also nicht der Einheitskreis, sondern ein wachsender Pfad in der komplexen Ebene.


c) \textbf{tanh ist sinnvoller.} Trade-offs:

\begin{enumerate}
  \item \textbf{Wertebereich}: tanh: ℝ → (−1, 1), nullzentriert und beschränkt → stabile Gradienten und gute Skalierung. sin: ℝ → [−1, 1], aber \textbf{periodisch}: zwei sehr verschiedene Eingaben (z. B. z=0 und z=2π) liefern denselben Output → mehrdeutige Repräsentation, schwierige Optimierung.
  \item \textbf{Monotonie}: tanh ist streng monoton ⇒ eindeutige Umkehrung möglich, kontinuierliches Lernsignal. sin oszilliert und hat unendlich viele lokale Extrema mit Ableitung 0 ⇒ Backpropagation würde durch periodisch verschwindende Gradienten gestört.
  \item (Bonus) \textbf{Sättigung}: tanh sättigt zwar weit weg von 0 (Vanishing Gradient), aber nur einmal pro Richtung; sin sättigt \textbf{immer wieder}, was Optimierer praktisch unkontrollierbar macht.
\end{enumerate}




% ──────────────────────────────────────────────────────────────

\end{document}
